% Options for packages loaded elsewhere
\PassOptionsToPackage{unicode}{hyperref}
\PassOptionsToPackage{hyphens}{url}
\documentclass[
]{article}
\usepackage{xcolor}
\usepackage[margin=1in]{geometry}
\usepackage{amsmath,amssymb}
\setcounter{secnumdepth}{-\maxdimen} % remove section numbering
\usepackage{iftex}
\ifPDFTeX
  \usepackage[T1]{fontenc}
  \usepackage[utf8]{inputenc}
  \usepackage{textcomp} % provide euro and other symbols
\else % if luatex or xetex
  \usepackage{unicode-math} % this also loads fontspec
  \defaultfontfeatures{Scale=MatchLowercase}
  \defaultfontfeatures[\rmfamily]{Ligatures=TeX,Scale=1}
\fi
\usepackage{lmodern}
\ifPDFTeX\else
  % xetex/luatex font selection
\fi
% Use upquote if available, for straight quotes in verbatim environments
\IfFileExists{upquote.sty}{\usepackage{upquote}}{}
\IfFileExists{microtype.sty}{% use microtype if available
  \usepackage[]{microtype}
  \UseMicrotypeSet[protrusion]{basicmath} % disable protrusion for tt fonts
}{}
\makeatletter
\@ifundefined{KOMAClassName}{% if non-KOMA class
  \IfFileExists{parskip.sty}{%
    \usepackage{parskip}
  }{% else
    \setlength{\parindent}{0pt}
    \setlength{\parskip}{6pt plus 2pt minus 1pt}}
}{% if KOMA class
  \KOMAoptions{parskip=half}}
\makeatother
\usepackage{color}
\usepackage{fancyvrb}
\newcommand{\VerbBar}{|}
\newcommand{\VERB}{\Verb[commandchars=\\\{\}]}
\DefineVerbatimEnvironment{Highlighting}{Verbatim}{commandchars=\\\{\}}
% Add ',fontsize=\small' for more characters per line
\usepackage{framed}
\definecolor{shadecolor}{RGB}{248,248,248}
\newenvironment{Shaded}{\begin{snugshade}}{\end{snugshade}}
\newcommand{\AlertTok}[1]{\textcolor[rgb]{0.94,0.16,0.16}{#1}}
\newcommand{\AnnotationTok}[1]{\textcolor[rgb]{0.56,0.35,0.01}{\textbf{\textit{#1}}}}
\newcommand{\AttributeTok}[1]{\textcolor[rgb]{0.13,0.29,0.53}{#1}}
\newcommand{\BaseNTok}[1]{\textcolor[rgb]{0.00,0.00,0.81}{#1}}
\newcommand{\BuiltInTok}[1]{#1}
\newcommand{\CharTok}[1]{\textcolor[rgb]{0.31,0.60,0.02}{#1}}
\newcommand{\CommentTok}[1]{\textcolor[rgb]{0.56,0.35,0.01}{\textit{#1}}}
\newcommand{\CommentVarTok}[1]{\textcolor[rgb]{0.56,0.35,0.01}{\textbf{\textit{#1}}}}
\newcommand{\ConstantTok}[1]{\textcolor[rgb]{0.56,0.35,0.01}{#1}}
\newcommand{\ControlFlowTok}[1]{\textcolor[rgb]{0.13,0.29,0.53}{\textbf{#1}}}
\newcommand{\DataTypeTok}[1]{\textcolor[rgb]{0.13,0.29,0.53}{#1}}
\newcommand{\DecValTok}[1]{\textcolor[rgb]{0.00,0.00,0.81}{#1}}
\newcommand{\DocumentationTok}[1]{\textcolor[rgb]{0.56,0.35,0.01}{\textbf{\textit{#1}}}}
\newcommand{\ErrorTok}[1]{\textcolor[rgb]{0.64,0.00,0.00}{\textbf{#1}}}
\newcommand{\ExtensionTok}[1]{#1}
\newcommand{\FloatTok}[1]{\textcolor[rgb]{0.00,0.00,0.81}{#1}}
\newcommand{\FunctionTok}[1]{\textcolor[rgb]{0.13,0.29,0.53}{\textbf{#1}}}
\newcommand{\ImportTok}[1]{#1}
\newcommand{\InformationTok}[1]{\textcolor[rgb]{0.56,0.35,0.01}{\textbf{\textit{#1}}}}
\newcommand{\KeywordTok}[1]{\textcolor[rgb]{0.13,0.29,0.53}{\textbf{#1}}}
\newcommand{\NormalTok}[1]{#1}
\newcommand{\OperatorTok}[1]{\textcolor[rgb]{0.81,0.36,0.00}{\textbf{#1}}}
\newcommand{\OtherTok}[1]{\textcolor[rgb]{0.56,0.35,0.01}{#1}}
\newcommand{\PreprocessorTok}[1]{\textcolor[rgb]{0.56,0.35,0.01}{\textit{#1}}}
\newcommand{\RegionMarkerTok}[1]{#1}
\newcommand{\SpecialCharTok}[1]{\textcolor[rgb]{0.81,0.36,0.00}{\textbf{#1}}}
\newcommand{\SpecialStringTok}[1]{\textcolor[rgb]{0.31,0.60,0.02}{#1}}
\newcommand{\StringTok}[1]{\textcolor[rgb]{0.31,0.60,0.02}{#1}}
\newcommand{\VariableTok}[1]{\textcolor[rgb]{0.00,0.00,0.00}{#1}}
\newcommand{\VerbatimStringTok}[1]{\textcolor[rgb]{0.31,0.60,0.02}{#1}}
\newcommand{\WarningTok}[1]{\textcolor[rgb]{0.56,0.35,0.01}{\textbf{\textit{#1}}}}
\usepackage{graphicx}
\makeatletter
\newsavebox\pandoc@box
\newcommand*\pandocbounded[1]{% scales image to fit in text height/width
  \sbox\pandoc@box{#1}%
  \Gscale@div\@tempa{\textheight}{\dimexpr\ht\pandoc@box+\dp\pandoc@box\relax}%
  \Gscale@div\@tempb{\linewidth}{\wd\pandoc@box}%
  \ifdim\@tempb\p@<\@tempa\p@\let\@tempa\@tempb\fi% select the smaller of both
  \ifdim\@tempa\p@<\p@\scalebox{\@tempa}{\usebox\pandoc@box}%
  \else\usebox{\pandoc@box}%
  \fi%
}
% Set default figure placement to htbp
\def\fps@figure{htbp}
\makeatother
\setlength{\emergencystretch}{3em} % prevent overfull lines
\providecommand{\tightlist}{%
  \setlength{\itemsep}{0pt}\setlength{\parskip}{0pt}}
\usepackage{bookmark}
\IfFileExists{xurl.sty}{\usepackage{xurl}}{} % add URL line breaks if available
\urlstyle{same}
\hypersetup{
  pdftitle={tidybayes Workflow},
  hidelinks,
  pdfcreator={LaTeX via pandoc}}

\title{tidybayes Workflow}
\author{}
\date{\vspace{-2.5em}2026-04-17}

\begin{document}
\maketitle

\subsection{Introduction}\label{introduction}

\texttt{tidybayes} is the bridge between Bayesian model objects and the
rest of the tidyverse. After fitting a model with \texttt{brms},
\texttt{rstanarm}, or raw Stan, you typically want to compare posteriors
across parameters, summarise interval estimates by subgroup, or layer
credible bands on top of raw observations. Doing this with the matrices
that samplers natively produce is awkward; \texttt{tidybayes} reshapes
them into long-format tibbles in which every draw of every parameter
occupies its own row, making the familiar \texttt{dplyr} verbs and
\texttt{ggplot2} geoms work without modification.

\subsection{Prerequisites}\label{prerequisites}

A working understanding of Bayesian inference (priors, likelihoods,
posteriors), comfort with \texttt{dplyr} pipelines, and a fitted MCMC
model object. You do not need to write Stan by hand --- \texttt{brm()}
or \texttt{stan\_glm()} is enough.

\subsection{Theory}\label{theory}

A posterior is a joint distribution over parameters represented as \(S\)
MCMC draws. \texttt{tidybayes::spread\_draws()} returns one row per draw
per parameter, so a model with \(K\) regression coefficients yields
\(S \times K\) rows. \texttt{gather\_draws()} produces an even longer
format with a \texttt{.variable} column, useful when faceting. The
package adds geoms --- \texttt{stat\_halfeye},
\texttt{stat\_pointinterval}, \texttt{stat\_lineribbon} --- that compute
density and interval summaries on the fly, so the data frame stays
untransformed. Predictive draws come from \texttt{epred\_draws()}
(expectation of the posterior predictive) and
\texttt{predicted\_draws()} (full predictive distribution including
residual variation).

\subsection{Assumptions}\label{assumptions}

The model has already converged: chains have mixed, \(\hat R\) is below
1.01, and effective sample size is adequate. \texttt{tidybayes} does not
diagnose convergence; it visualises whatever draws you give it, so a
poorly mixed chain produces a misleading posterior summary.

\subsection{R Implementation}\label{r-implementation}

\begin{Shaded}
\begin{Highlighting}[]
\FunctionTok{library}\NormalTok{(brms); }\FunctionTok{library}\NormalTok{(tidybayes); }\FunctionTok{library}\NormalTok{(ggplot2)}

\FunctionTok{data}\NormalTok{(mtcars)}
\NormalTok{fit }\OtherTok{\textless{}{-}} \FunctionTok{brm}\NormalTok{(mpg }\SpecialCharTok{\textasciitilde{}}\NormalTok{ wt }\SpecialCharTok{+}\NormalTok{ hp, }\AttributeTok{data =}\NormalTok{ mtcars, }\AttributeTok{chains =} \DecValTok{4}\NormalTok{, }\AttributeTok{iter =} \DecValTok{2000}\NormalTok{, }\AttributeTok{refresh =} \DecValTok{0}\NormalTok{)}

\CommentTok{\# Tidy draws}
\NormalTok{draws }\OtherTok{\textless{}{-}}\NormalTok{ fit }\SpecialCharTok{|\textgreater{}}
  \FunctionTok{spread\_draws}\NormalTok{(b\_wt, b\_hp)}

\FunctionTok{head}\NormalTok{(draws)}

\CommentTok{\# Visualise posterior}
\NormalTok{draws }\SpecialCharTok{|\textgreater{}}
  \FunctionTok{pivot\_longer}\NormalTok{(}\FunctionTok{c}\NormalTok{(b\_wt, b\_hp), }\AttributeTok{names\_to =} \StringTok{"parameter"}\NormalTok{, }\AttributeTok{values\_to =} \StringTok{"value"}\NormalTok{) }\SpecialCharTok{|\textgreater{}}
  \FunctionTok{ggplot}\NormalTok{(}\FunctionTok{aes}\NormalTok{(}\AttributeTok{x =}\NormalTok{ value, }\AttributeTok{y =}\NormalTok{ parameter)) }\SpecialCharTok{+}
  \FunctionTok{stat\_halfeye}\NormalTok{(}\AttributeTok{fill =} \StringTok{"\#2A9D8F"}\NormalTok{) }\SpecialCharTok{+}
  \FunctionTok{theme\_minimal}\NormalTok{()}

\CommentTok{\# Conditional effects}
\NormalTok{fit }\SpecialCharTok{|\textgreater{}}
  \FunctionTok{epred\_draws}\NormalTok{(}\AttributeTok{newdata =} \FunctionTok{data.frame}\NormalTok{(}\AttributeTok{wt =} \FloatTok{2.5}\NormalTok{, }\AttributeTok{hp =} \FunctionTok{seq}\NormalTok{(}\DecValTok{50}\NormalTok{, }\DecValTok{300}\NormalTok{, }\DecValTok{10}\NormalTok{))) }\SpecialCharTok{|\textgreater{}}
  \FunctionTok{ggplot}\NormalTok{(}\FunctionTok{aes}\NormalTok{(}\AttributeTok{x =}\NormalTok{ hp, }\AttributeTok{y =}\NormalTok{ .epred)) }\SpecialCharTok{+}
  \FunctionTok{stat\_lineribbon}\NormalTok{(}\AttributeTok{alpha =} \FloatTok{0.3}\NormalTok{) }\SpecialCharTok{+}
  \FunctionTok{theme\_minimal}\NormalTok{()}
\end{Highlighting}
\end{Shaded}

\subsection{Output \& Results}\label{output-results}

\texttt{spread\_draws()} returns a tibble with \texttt{.chain},
\texttt{.iteration}, \texttt{.draw}, and one column per parameter. The
half-eye plot shows the posterior density together with a black point at
the median and bars at the 66 \% and 95 \% credible intervals.
\texttt{stat\_lineribbon} produces a shaded band that widens at
horsepower values poorly supported by the data, communicating where
extrapolation is risky.

\subsection{Interpretation}\label{interpretation}

A reporting sentence: ``The posterior median effect of weight on
miles-per-gallon was \(-3.9\) (95 \% CrI \(-5.4\) to \(-2.4\)), while
horsepower contributed \(-0.03\) (95 \% CrI \(-0.05\) to \(-0.01\));
both intervals exclude zero, consistent with each predictor having a
non-trivial association after the other is controlled for.'' Pair
narrative numbers with the half-eye plot so readers can see the full
posterior shape and not just an interval.

\subsection{Practical Tips}\label{practical-tips}

\begin{itemize}
\tightlist
\item
  Use \texttt{spread\_draws()} when you want one column per parameter
  and \texttt{gather\_draws()} when you plan to facet by parameter name.
\item
  Regular expressions are accepted:
  \texttt{spread\_draws(fit,\ b\_{[}a-z{]}+)} extracts every regression
  coefficient at once.
\item
  \texttt{add\_epred\_draws()} and \texttt{add\_predicted\_draws()}
  attach predictive draws directly to a \texttt{newdata} frame, which is
  convenient when overlaying credible bands on observed points.
\item
  Use \texttt{mean\_qi()}, \texttt{median\_qi()}, or
  \texttt{mode\_hdi()} to collapse draws into interval summaries;
  \texttt{point\_interval\ =\ mean\_qi} is the default in most geoms but
  is worth setting explicitly for reproducibility.
\item
  The package is engine-agnostic: the same code path supports
  \texttt{rstan}, \texttt{cmdstanr}, \texttt{rstanarm}, and
  \texttt{brms}, so refactoring between them rarely affects downstream
  plotting code.
\item
  Avoid coercing draws to wide format prematurely; \texttt{ggplot2} is
  happiest with the long tibble that \texttt{tidybayes} returns.
\end{itemize}

\subsection{R Packages Used}\label{r-packages-used}

\texttt{brms} for model fitting, \texttt{tidybayes} for posterior
reshaping and stat geoms, \texttt{ggplot2} and \texttt{dplyr} for
plotting and manipulation, plus \texttt{posterior} (loaded transitively)
for low-level draw arrays.

\subsection{For Reviewers}\label{for-reviewers}

\textbf{What to look for in a paper using this method.}

\begin{itemize}
\tightlist
\item
  \textbf{Common misapplications.}
\item
  \textbf{Diagnostics that should be reported but often aren't.}
\item
  \textbf{Red flags in tables and figures.}
\item
  \textbf{What to verify.}
\item
  \textbf{What an adequate Methods paragraph must contain.}
\end{itemize}

\subsection{See also --- labs in this
chapter}\label{see-also-labs-in-this-chapter}

\begin{itemize}
\tightlist
\item
  \href{labs/lab-bayesian-thinking-control-vs-decision-error.Rmd}{Bayesian
  thinking; α control vs decision error}
\item
  \href{labs/lab-brms-stan-loo-hierarchical-models.Rmd}{brms / Stan,
  LOO, hierarchical models}
\end{itemize}

\end{document}
